\chapter{Lesson 10: Adjectives, Likes and Dislikes}

\section{Two kinds of adjectives}

Japanese adjectives conjugate (for tense and negation) like verbs do. There are two
types:
\begin{itemize}
  \item \textbf{\ja{い}-adjectives} end in \ja{い} in their dictionary form:
        \ja{たかい} expensive, high; \ja{おもしろい} interesting.
  \item \textbf{\ja{な}-adjectives} take \ja{な} before a noun: \ja{しずかなまち}
        ``a quiet town''. They behave much like nouns in the predicate.
\end{itemize}

\subsection{\ja{い}-adjectives}
\begin{center}
\begin{tabular}{lll@{\qquad}lll}
\toprule
\ja{おおきい} & ōkii & big & \ja{ちいさい} & chiisai & small \\
\ja{たかい} & takai & expensive; tall & \ja{やすい} & yasui & cheap \\
\ja{あたらしい} & atarashii & new & \ja{ふるい} & furui & old (things) \\
\ja{あつい} & atsui & hot (weather) & \ja{さむい} & samui & cold (weather) \\
\ja{おもしろい} & omoshiroi & interesting & \ja{つまらない} & tsumaranai & boring \\
\ja{むずかしい} & muzukashii & difficult & \ja{やさしい} & yasashii & easy; kind \\
\ja{いそがしい} & isogashii & busy & \ja{こわい} & kowai & scary \\
\ja{おいしい} & oishii & delicious & \ja{いい} & ii & good \\
\bottomrule
\end{tabular}
\end{center}

\subsection{\ja{な}-adjectives}
\begin{center}
\begin{tabular}{lll@{\qquad}lll}
\toprule
\ja{げんき} & genki & healthy, energetic & \ja{しずか} & shizuka & quiet \\
\ja{にぎやか} & nigiyaka & lively & \ja{きれい} & kirei & beautiful, clean \\
\ja{ひま} & hima & not busy, free & \ja{ゆうめい} & yūmei & famous \\
\ja{すき} & suki & liked & \ja{きらい} & kirai & disliked \\
\ja{だいすき} & daisuki & very much liked & \ja{だいきらい} & daikirai & hated \\
\ja{しんせつ} & shinsetsu & kind & \ja{たいへん} & taihen & tough, hard \\
\bottomrule
\end{tabular}
\end{center}

\begin{important}[Traps]
\ja{きれい} and \ja{ゆうめい} end in the sound \emph{i} but are \ja{な}-adjectives
(the \ja{い} is part of a long vowel \emph{ē}). Similarly \ja{きらい} is a \ja{な}-adjective.
Their negatives are \ja{きれいじゃないです}, never *\ja{きれくないです}.
\end{important}

\section{Using adjectives}

\subsection{Before a noun}
\begin{itemize}
  \item \ja{い}-adjective: dictionary form + noun: \ja{おもしろいほん} an interesting book.
  \item \ja{な}-adjective: adjective + \ja{な} + noun: \ja{しずかなへや} a quiet room.
\end{itemize}

\subsection{As a predicate (``X is Adj'')}
\ja{このほんはおもしろいです。} This book is interesting.\\
\ja{このへやはしずかです。} This room is quiet.

\section{Conjugation}

\begin{center}
\begin{tikzpicture}[x=1cm,y=1cm, c/.style={draw, rounded corners, minimum width=3.5cm, minimum height=0.75cm, font=\small}]
  \node[font=\bfseries\small] at (0,1.1) {\ja{い}-adjective \ja{たかい}};
  \node[c, fill=blue!6] (p) at (0,0) {\ja{たか\textbf{い}です}};
  \node[c, fill=red!6] (n) at (0,-1.1) {\ja{たか\textbf{くない}です}};
  \node[c, fill=blue!6] (pp) at (0,-2.2) {\ja{たか\textbf{かった}です}};
  \node[c, fill=red!6] (np) at (0,-3.3) {\ja{たか\textbf{くなかった}です}};
  \node[font=\bfseries\small] at (6,1.1) {\ja{な}-adjective \ja{しずか}};
  \node[c, fill=blue!6] (p2) at (6,0) {\ja{しずか\textbf{です}}};
  \node[c, fill=red!6] (n2) at (6,-1.1) {\ja{しずか\textbf{じゃない}です}};
  \node[c, fill=blue!6] (pp2) at (6,-2.2) {\ja{しずか\textbf{でした}}};
  \node[c, fill=red!6] (np2) at (6,-3.3) {\ja{しずか\textbf{じゃなかった}です}};
  \node[font=\small] at (3,0) {present aff.};
  \node[font=\small] at (3,-1.1) {present neg.};
  \node[font=\small] at (3,-2.2) {past aff.};
  \node[font=\small] at (3,-3.3) {past neg.};
\end{tikzpicture}
\end{center}

\begin{itemize}
  \item \ja{い}-adjectives: drop the final \ja{い}; add \ja{くない} (negative),
        \ja{かった} (past), \ja{くなかった} (past negative); then \ja{です} for
        politeness. \textbf{Never} use \ja{でした} with an \ja{い}-adjective:
        *\ja{たかいでした} is wrong.
  \item \ja{な}-adjectives conjugate exactly like nouns + \ja{です} (Lesson 9).
\end{itemize}

\begin{important}[The irregular \ja{いい}]
\ja{いい} ``good'' conjugates from its older form \ja{よい}:
\ja{いいです}, \ja{よくないです}, \ja{よかったです}, \ja{よくなかったです}.
Same for \ja{かっこいい} ``cool'': \ja{かっこよかったです}.
\end{important}

\begin{example}
\leavevmode
\begin{itemize}
  \item \ja{きのうはさむかったです。} \emph{Kinō wa samukatta desu.} It was cold yesterday.
  \item \ja{テストはむずかしくなかったです。} The test wasn't difficult.
  \item \ja{りょこうはたのしかったです。} The trip was fun.
  \item \ja{そのまちはにぎやかじゃなかったです。} That town wasn't lively.
  \item \ja{えいがはよくなかったです。} The movie wasn't good.
\end{itemize}
\end{example}

\subsection{Degree adverbs}
\ja{とても} very, \ja{ちょっと} a little; \ja{あまり} (not) very and \ja{ぜんぜん}
(not) at all with negatives:
\ja{このコーヒーはあまりおいしくないです。} This coffee isn't very good.

\section{Likes and dislikes: \ja{すき}/\ja{きらい}}

\ja{すき} and \ja{きらい} are \ja{な}-adjectives. The thing liked is marked with \ja{が}
(not \ja{を}!).
\begin{center}
\textbf{Person \ja{は} thing \ja{が} \ja{すき}/\ja{きらい}です}
\end{center}

\begin{example}
\leavevmode
\begin{itemize}
  \item \ja{わたしはねこがすきです。} I like cats.
  \item \ja{やさいがあまりすきじゃないです。} I don't like vegetables very much.
  \item \ja{なにがすきですか。} What do you like?
  \item \ja{どんなおんがくがすきですか。} What kind of music do you like?
        --- \ja{ジャズがだいすきです。} I love jazz.
  \item \ja{すきなたべものはなんですか。} What is your favourite food? (\ja{すきな} + noun)
\end{itemize}
\end{example}

\begin{aside}
\ja{きらい} is a strong word. Saying \ja{きらいです} about someone's cooking or hobby is
blunt; Japanese speakers usually soften dislikes to \ja{あまりすきじゃないです}
(``I'm not very fond of it''). Also \ja{すきです} said directly to a person is a romantic
confession, not ``I like you as a friend''.
\end{aside}

\section{Other useful patterns}

\subsection{\ja{どんな} ``what kind of''}
\ja{どんなひとですか。} ``What kind of person is (s)he?'' ---
\ja{やさしいひとです。} ``(S)he is a kind person.''

\subsection{\ja{ね} and \ja{よ}}
\begin{itemize}
  \item \ja{ね} seeks agreement / shows shared feeling: \ja{いいてんきですね。}
        ``Nice weather, isn't it?''
  \item \ja{よ} gives information the listener doesn't have: \ja{あのレストランは
        たかいですよ。} ``That restaurant is expensive, you know.''
\end{itemize}

\subsection{Counting people}
\ja{ひとり} 1 person, \ja{ふたり} 2 people, then \ja{さんにん}, \ja{よにん},
\ja{ごにん}, \ldots ``How many people?'' \ja{なんにん}.

\section{Kanji}
\begin{center}
\begin{tabular}{lllll}
\toprule
Kanji & Meaning & Readings & Example & \\
\midrule
\ja{大} & big & \ja{ダイ}, \ja{タイ}; \ja{おお(きい)} & \ja{大学} & daigaku, university \\
\ja{小} & small & \ja{ショウ}; \ja{ちい(さい)} & \ja{小さい} & chiisai, small \\
\ja{高} & high, expensive & \ja{コウ}; \ja{たか(い)} & \ja{高校} & kōkō, high school \\
\ja{安} & cheap, peaceful & \ja{アン}; \ja{やす(い)} & \ja{安い} & yasui, cheap \\
\ja{新} & new & \ja{シン}; \ja{あたら(しい)} & \ja{新聞} & shinbun, newspaper \\
\ja{古} & old & \ja{コ}; \ja{ふる(い)} & \ja{古い} & furui, old \\
\ja{好} & like & \ja{コウ}; \ja{す(き)} & \ja{好き} & suki, liked \\
\bottomrule
\end{tabular}
\end{center}
Kun readings show the okurigana in parentheses: the part written in hiragana after the
kanji.

\begin{center}
\begin{tikzpicture}[k/.style={draw, minimum size=1.2cm, font=\huge}]
  \node[k] (a) at (0,0) {\ja{女}};
  \node[font=\Large] at (1.1,0) {+};
  \node[k] (b) at (2.2,0) {\ja{子}};
  \node[font=\Large] at (3.3,0) {=};
  \node[k, fill=blue!8] (c) at (4.4,0) {\ja{好}};
  \node[font=\small] at (0,-0.95) {woman};
  \node[font=\small] at (2.2,-0.95) {child};
  \node[font=\small] at (4.4,-0.95) {like, fond of};
\end{tikzpicture}
\end{center}

\section*{Lesson 10 review}
\begin{itemize}
  \item \ja{い}-adj: \ja{〜い}, \ja{〜くない}, \ja{〜かった}, \ja{〜くなかった} (+ \ja{です}).
  \item \ja{な}-adj: like nouns: \ja{です}, \ja{じゃない}, \ja{でした}, \ja{じゃなかった}.
  \item \ja{いい} $\to$ \ja{よくない}, \ja{よかった}.
  \item Before nouns: \ja{たかいほん}, \ja{しずかなへや}.
  \item X \ja{が} \ja{すき}/\ja{きらい}です; soften dislikes.
  \item Kanji: \ja{大小高安新古好}.
\end{itemize}
